\documentclass[10pt]{article}
\usepackage[a4paper,margin=0.8in]{geometry}
\usepackage{amsmath,amssymb}
\usepackage{hyperref}
\title{Six central automorphisms: a counterexample to Conjecture 00000004274}
\author{Local solution draft}
\date{October 3, 2026}
% Compact consistent title for a short mathematical note.
\makeatletter
\renewcommand{\maketitle}{\begin{center}
{\Large\bfseries\@title\par}\vspace{0.6em}
{\small\@author\par\@date}
\end{center}\vspace{0.5em}}
\makeatother
\setlength{\emergencystretch}{3em}
\begin{document}
\maketitle
\noindent\textbf{Verdict: false under the standard integer-power counting interpretation.} The automorphism group of the countable $7$-group $\mathbb Z/7\mathbb Z$ is abelian of order $6$, which is not a power of $2$.

\paragraph{Statement being refuted.}
Conjecture 00000004274 states that the number of central automorphisms of a countable $p$-group is given by a termwise product of powers of $2$ read from its Ulm sequence. In a finite counting formula, the exponents are nonnegative integers (equivalently finite cardinal exponents). A finite product of such powers is again a power of $2$:
\[
 \prod_{j=1}^{r}2^{e_j}=2^{\sum_{j=1}^{r}e_j}.
\]
We refute this necessary consequence, independently of the other assertion involving ``characteristic elements.''

\paragraph{The group and all its automorphisms.}
Take the additive group $G=\mathbb Z/7\mathbb Z$. It is finite, hence countable, and every element is killed by $7$. Every additive homomorphism $f:G\to G$ satisfies
\[
 f(x)=x f(1).
\]
Indeed, each element is a sum of copies of $1$, and additivity preserves those sums. Such a homomorphism is an automorphism exactly when $f(1)\ne0$. Thus all automorphisms are precisely
\[
 m_a(x)=ax,\qquad a\in\{1,2,3,4,5,6\}\pmod7.
\]
They are pairwise distinct by evaluation at $1$. Their inverse multipliers, in this order, are
\[
 1,4,5,2,3,6.
\]
For every $a,b$, composition gives $m_a m_b=m_{ab}=m_{ba}=m_bm_a$. Therefore every automorphism is central and
\[
 \big|Z(\operatorname{Aut}(G))\big|=6.
\]
The powers $2^0,2^1,2^2$ are $1,2,4$, and all later powers are divisible by $8$. None is $6$. Consequently the required power-of-two product cannot count this center.

The group has only one nonzero Ulm invariant, $u_0=1$, since $7G=0$. Hence there is no issue involving a genuinely infinite nontrivial product. Even if one retains the trailing zero invariants, inserting factors equal to $1$ cannot change the obstruction.

\paragraph{Formal verification and semantic correspondence.}
The accompanying \texttt{Main.lean} imports only \texttt{Std} and is checked with Lean 4.19.0. The carrier is \texttt{Fin 7}, with modular addition and multiplication. The file checks the abelian group laws, primality of $7$, countability of the carrier, and annihilation by $7$.

An automorphism is an actual function on this carrier with proofs of additivity, injectivity, and surjectivity. The proof that every homomorphism is determined by $f(1)$ is symbolic induction, not a test over a selected family of maps. Six multiplier automorphisms are constructed, and inverse maps between the full automorphism type and the six-element index set are proved. Centrality is defined by commuting with \emph{every} automorphism as functions, and that property is proved for all of them.

The actual center is then enumerated. Completeness and absence of duplicates are proved before its list length is computed as $6$. The final theorem refutes the claimed product formula for a complete, duplicate-free enumeration of this center. Thus the structural bridge from the group to the count is part of the checked argument; the final theorem is not an isolated numerical assertion about $6$.

The axiom audit reports only the standard Lean axioms \texttt{propext} and \texttt{Quot.sound}. No extra axioms, admissions, or native decision procedures occur.

\paragraph{Interpretation boundary.}
The source does not explicitly give a domain for the exponents in ``powers of $2$.'' The counterexample applies to the usual nonnegative-integer or cardinal exponents of a counting formula. Allowing arbitrary real exponents would make every positive finite count a power of $2$, via $N=2^{\log_2 N}$, and would remove the claimed counting restriction altogether. This is a semantic clarification required by the source, not a computational uncertainty.
\end{document}
